
\section{Contenido de TFG y TFM} 

\subsubsection{Estructura y contenidos de un Trabajo de Investigación}
\label{sec:estructuratrabajoinvestigacion}

Un TFG o TFM en la modalidad de Trabajo de Investigación deberá estructurarse en los siguientes apartados:
\begin{enumerate}
    \item \textit{Introducción}: contendrá un estudio del estado actual de la técnica sobre el problema plateado, las alternativas consideradas y los objetivos a conseguir.
    \item \textit{Material y métodos}: se describirán los equipos y materiales utilizados para su realización, la metodología empleada y una descripción completa de los experimentos realizados.
    \item \textit{Resultados y discusión}: se describirán y analizarán los resultados obtenidos siguiendo la metodología propuesta.
    \item \textit{Conclusiones}: en función de los resultados obtenidos, 
    se expondrán las conclusiones de la investigación y se propondrán trabajos futuros.
    \item \textit{Anexos}: tantos como se considere necesario en función de los aspectos concretos que se quieran especificar.
    \item \textit{Bibliografía}: se indicarán todas las referencias bibliográficas utilizadas, ordenadas con código para su cita en la memoria.
\end{enumerate}